% !TeX root = ../papers/nc-p3.tex
\section{Bell曲面の全流入と十五次元成分およびその滑らかな近傍}
\label{sec:split-surface-component}
\begin{theorem}\label{thm:split-surface-component}
$u v(u^2-1)(v^2-1)\ne0$ における前節の四つの明示的な二平面分解族は
全て定理\ref{thm:full-fifteen-component}の十五次元成分に属する。
各族の一般の閉軌道と、それに至る\emph{全て}の無印の窓で、
隣接二切断の階数スキームはアフィン錐上で滑らかな76次元であり、この成分だけが通る。
さらにそれらの全点を含む不変なザリスキー開集合で、全窓は標準Koszul左右AS正則で、
四次中心は五次元、代数は四つの中心四次式の生成する部分環上有限である。
一般閉軌道での射影切片の接表現のトーラス重み $-2,-1,0,1,2$ の多重度は $3,4,2,4,3$ である。
\end{theorem}
\begin{proof}[五中心の計算と量子特殊化]
四位置で二進表示した四文字の排他的論理和により各関係に四元群次数を与える。
両行列切断の六関係を $\mathbb Q(u,v)$ で取り、各始点で経路長5まで一次関係を挿入して
経路基底を取り、一文字との可換式の両側を射影する。
長さ4の中心切断の次元は四文字次数について $(5,0,0,0)$ である。
\path{computations/verify_split_surface_component.py} は四つの $h$ 全てで、
有理関数体上の経路基底から計224行・44列となる自明指標の可換式および残り三指標を計算する。

$u=\iota$, $v=\iota$ を $h=X,XZ$ に、$u=\iota$, $v=-\iota$ を $h=1,Z$ に使う。
全て $\iota^2=-1$ とする。どちらの四文字表示も正則で、二切断の双方の
相補逆行列の分母は単元である。前者は位置2,3、後者は位置1,3で各二文字のビットを反転する。
非零の24語は全て四文字の置換になる。各辺の二つの文字 $a,b$ の交換係数を $q_{i,ab}$ とすると、
前記の検算で、各向きの相補二文字をどちらの順序にとっても同じ比が得られる。
四つの三角積 $q_{i,ab}q_{i,bc}/q_{i,ac}$ ($a<b<c$) は全 $i$ で同一の
$(\iota,-\iota,\iota,-\iota)$ またはその逆数である。
従って相続く位置の交換係数の比は文字ごとの倍率の比であり、添字付きの対角基底変更で
一定な量子交換係数にできる。この三角積は定理\ref{thm:quantum-component-intersection}の
唯一の位数4の量子類を与える。

二変数族の全点で経路の各次数は標準次元なので、四次中心の可換式は束の射であり、
五次元以上の条件は閉である。四族とも既約で、一般点で五次元なので全体がこの閉条件に入る。
量子点の近傍では定理\ref{thm:quantum-ore-incidence}により、この閉条件は
既知の十五次元成分そのものである。上の特殊化を用いると、各既約族全体が同じ成分に入る。
\end{proof}
\begin{proof}[全接方向、全流入と正則開集合]
$u=2,v=3$ を素数109で還元し、定理\ref{thm:isolated-split-component}の
左右核の全200行・256列で接方程式を取る。法線数 $0,1,2,3,4$ のアフィン接次元は
四補完とも $3,20,30,20,3$ である。各切断の階数6は保たれており、
標数零の幾何学的一般点でもこの和76が接次元の上界である。
既知成分のアフィン次元が76だから等号で、局所環は正則、通る成分は一つだけである。
射影化と軌道方向を除くには、重み $\pm1$ のせん断各16、重み零の28
（ブロック対角群の次元32から三つの自明スカラーと一つのトーラスを除く）を引く。
この28にはアフィンの共通倍率の方向も入る。これで表示した多重度を得る。
滑らかでこの一成分だけに属する条件は基底群不変な開条件なので、
それを満たす点を軌道閉包に持つ全窓にも成立する。

最後に量子点では通常の四乗四元が中心で、これらの商は
各文字の指数が3以下の量子単項式で張られる。標準四移動表示への輸送では中心切断として取る。
量子点の四次中心は五次元であり、この成分上で一般点も五次元なので量子点の近傍で束となる。
その四元を局所切断に持ち上げ、中心商の有限性と標準解消を認証する有限次数検査の
開性（定理\ref{thm:central-window-cutoff}）を適用する。
これにより量子点を含む成分内の開集合の全点で主張が成立する。
その基底群による全ての移動の和も同様で、四族の各一般点は上の特殊化からこの開集合に入る。
さらに各一般点で得た滑らかな一成分だけの近傍に制限すれば、全階数スキーム内の不変開集合が得られる。
軌道閉包にその一般点を持つ全窓もこの開集合に入る。
\end{proof}
